From - Thu Jun  1 15:32:24 2000
Received: from mr.usu.edu (mr.usu.edu [129.123.1.3])
	by math.usu.edu (8.9.3+Sun/8.9.1) with ESMTP id WAA20342
	for <symanzik@sunfs.math.usu.edu>; Tue, 30 May 2000 22:32:58 -0600 (MDT)
Received: from hotmail.com (panda.uwrl.usu.edu [129.123.9.28])
	by mr.usu.edu (Postfix) with ESMTP id 5E85C1D2BC
	for <symanzik@sunfs.math.usu.edu>; Tue, 30 May 2000 22:34:50 -0600 (MDT)
Message-ID: <393497F9.5CF886D8@hotmail.com>
Date: Tue, 30 May 2000 22:41:29 -0600
From: Chian <chian98@hotmail.com>
X-Mailer: Mozilla 4.61 [en] (Win98; U)
X-Accept-Language: en
MIME-Version: 1.0
To: symanzik@sunfs.math.usu.edu
Subject: HW#9
Content-Type: multipart/mixed;
 boundary="------------B413C6379C9F3B1BE937284A"
Content-Length: 3103
Status: RO
X-Mozilla-Status: 8001
X-Mozilla-Status2: 00000000
X-UIDL: 38baa60a0000035c

This is a multi-part message in MIME format.
--------------B413C6379C9F3B1BE937284A
Content-Type: text/plain; charset=us-ascii
Content-Transfer-Encoding: 7bit

It's seems I forget to attach the HW#9 in the previous email. I resend
it. Thanks.

Qian

--------------B413C6379C9F3B1BE937284A
Content-Type: text/plain; charset=us-ascii;
 name="Hw"
Content-Transfer-Encoding: 7bit
Content-Disposition: inline;
 filename="Hw"

Stat5810      Hw#9 PartII      05/28/00
Zhao Qian


Evaluations of fellow students.

#2 presentation: Statlets,
http://www.statlets.com/statletsindex.htm 

In general, I understand how to use the statlets package. There will appear a window 
where we can input the data. It is very easy to input data using 'copy' and 'paste' 
technique. Then we can analyze the data using the functions shown on the tool bar. It 
is quite straight forward in terms of input data and analyze it. There are colorful plot 
as output. We can do the common statistical analysis and plot on the data such as: a). plot:
scatter plot, histogram, boxplot, bar chart, etc; b). analysis: summery statistics, T-test, 
simple regression, multiple regression, etc. However, there are several disadvantages in
this web-based package. The homepage is not very clear on how we should start to use this 
package. If not for the help of the instructor, it will need a lot of time to find out 
where we can start. Compare to Webstat, the method of input data is not complete and more 
complex. We do have enough time to follow the instructor but there are a couple of times 
that we are confused which botton to click. Since this statistic software is quite complete,
it is impossible to cover all the content of this package. However, through the exercise 
in class we are able to do some simple analysis with this package. 

#3 presentation: XploRe
http://www.xplore-stat.de  

This package is obviously more difficult for a beginner. The hompage of Xplore is not easy 
to understand. We need to download this software before we can use it. The software has 
already be download so we don't need to do this in class. But when I work on my own machine,
I have to fill out the form and download it. In general, this software is very good. It is 
easy to input the data using 'copy' and 'paste'. To do the analysis, one need to write a 
program using the language similar to S/S+. Follow the produres the instructor given, we 
are able to do some simple analysis such as finding the statistical summary, making  histogram
and boxplot. I think this package is somewhat like SAS. Compare to the previous mentioned 
web-based package such as Statlets and Webstat, this package is not as easy to understand,
So for a beginner in statistics, it is not very inviting. However, I think it is more powerful
than those two. It is able to do much more complex analysis. While other two are better in 
teaching the beginner. Since the limit of time, only very small part of the package is tought.
Also, the meaning of each line in the program is not explained clearly by the instructor.

--------------B413C6379C9F3B1BE937284A--

